\chapter{Antenna Patterns, Polarization, and Practical Design}
\label{ch:antennapatterns}
An antenna is an impedance at its feed point and a field pattern in space.  The
first view belongs to \cref{ch:lines}; this chapter develops the second.  The
same phasor addition that made the two-element array factor in
\cref{sec:arrayfactor} also explains beamwidth, front-to-back ratio,
polarization, and the behavior of practical directional antennas.

\section{A Pattern Is a Normalized Power Function}
\label{sec:patternmetrics}
At a fixed distance in the far field, let the complex electric-field amplitude
in direction \((\theta,\phi)\) be \(E(\theta,\phi)\).  Power density is
proportional to \(|E|^2\), so a useful dimensionless pattern is
\[
p(\theta,\phi)=\frac{|E(\theta,\phi)|^2}
                       {|E|^2_{\max}}.
\]
Its peak is one; it says where the antenna directs the power without hiding the
answer in a choice of transmitter power or distance.  A polar plot usually shows
\(10\log_{10}p\) in decibels.  Thus its \(-\SI{3}{\decibel}\) contour is the
\textbf{half-power beamwidth}: the angular width between directions for which
\(p=1/2\).  It is not an arbitrary convention: \(-\SI{3.01}{\decibel}\) is
exactly the logarithm of one half (\cref{sec:decibels}).

The \textbf{front-to-back ratio} makes the same comparison in two particular
directions:
\[
\boxed{\quad\mathrm{F/B}=10\log_{10}\frac{p_{\rm front}}{p_{\rm back}}\quad}.
\]
A \SI{20}{\decibel} ratio means the chosen rear direction receives one hundredth
as much power as the main direction.  It does \emph{not} state the gain: a pattern
can have a sharp main lobe and substantial loss, or modest directivity and high
efficiency.  Gain combines directivity with the radiation efficiency from
\cref{sec:antennaeff}.

\begin{exambox}
Beamwidth, gain, and front-to-back ratio are power statements.  Field amplitude
uses \(20\log_{10}\); pattern power uses \(10\log_{10}\).  The two conventions
give the same number only because power is proportional to amplitude squared.
\end{exambox}

\section{Polarization Is the Path Traced by the Field Vector}
\label{sec:polarization}
At one point in space, resolve a transverse wave's electric field into two
orthogonal components:
\[
 E_x=E_{0x}\cos\omega t,\qquad
 E_y=E_{0y}\cos(\omega t+\delta).
\]
The tip of \(\mathbf E\) traces a line when the components have phase difference
\(\delta=0\) or \(\pi\): this is \textbf{linear polarization}.  With equal
amplitudes and \(\delta=\pm\pi/2\), it traces a circle, giving right- or
left-hand \textbf{circular polarization}; unequal amplitudes give an ellipse.
The word ``hand'' is a viewing convention, so diagrams and specifications must
state the direction of propagation.

A receiving antenna responds to the component of \(\mathbf E\) aligned with it.
For two linearly polarized antennas with angle \(\alpha\) between their axes,
the received field is proportional to \(\cos\alpha\), hence the ideal received
power is
\[
\boxed{\quad P_r\propto\cos^2\alpha.\quad}
\]
Crossing vertical and horizontal Yagis therefore permits either linear
polarization; driving equal currents in quadrature makes circular polarization.
The calculation also explains why a perfect cross-polarized linear antenna has
an ideal null.  Real antennas, reflections, and imperfect alignment fill it in.

\section{Why Apertures Gain \SI{6}{\decibel} When Frequency Doubles}
\label{sec:aperturegain}
A reflector or dish collects a portion of a passing wave over its physical
aperture \(A\).  Its directivity must compare that area with the wavelength-sized
area occupied by one spatial mode, so
\[
\boxed{\quad G\approx\eta_{\rm ap}\frac{4\pi A}{\lambda^2},\qquad
0<\eta_{\rm ap}<1.\quad}
\]
Here \(\eta_{\rm ap}\) is aperture efficiency: illumination taper, blockage,
surface error, and spillover prevent a real dish from using every square metre.
For a fixed dish, \(G\propto1/\lambda^2\propto f^2\).  Doubling frequency
quadruples gain, which is \(10\log_{10}4=\SI{6.02}{\decibel}\).  Conversely,
keeping gain while doubling wavelength requires four times the aperture area.

This is a spatial version of the scaling already familiar from transmission
lines: wavelength determines how large an object looks electrically.  A dish
that is many wavelengths across can align many phase contributions into a narrow
lobe; a dish that is only a fraction of a wavelength across cannot.

\section{From Ideal Wires to Buildable Antennas}
The half-wave dipole and ideal ground image are powerful first models, not a
guarantee that a particular installation has their impedance or pattern.  A
thin-wire numerical model makes the approximation explicit: divide each wire
into short segments, assign an unknown current to each, and impose the boundary
condition that the conductor's tangential electric field vanishes.  The result is
a linear system
\[
\mathbf Z\mathbf I=\mathbf V.
\]
The entries of \(\mathbf Z\) represent mutual coupling between segments.  Once
the currents are known, their re-radiated fields are summed to estimate feed-point
impedance and pattern.  This is the \textbf{method of moments} in its simplest
form.  It is the antenna analogue of replacing a distributed line by many small
\(R'\), \(L'\), \(G'\), and \(C'\) sections: more, shorter segments generally
give a more faithful model, but measurements still decide the final installation.

Three design consequences follow directly from current distribution.
\begin{description}
\item[Loading is not free length.] A short vertical has little radiation
resistance and large capacitive reactance.  A loading coil cancels reactance at one
frequency, but coil and ground losses remain in series with the small radiation
resistance.  Its efficiency is therefore still the ratio in \cref{sec:antennaeff}.
\item[Put reactance where current is useful.] A centre or upper loading coil lies
where more wire current can contribute to radiation than a base coil.  Top loading
also lengthens the effective current path, reducing the inductance required.  The
precise optimum depends on the physical antenna and ground system; the current
distribution tells why these choices help.
\item[Parasitic elements are coupled resonators.] A Yagi's driven element excites
nearby detuned conductors.  A slightly longer reflector and slightly shorter director
re-radiate with different phases.  Spacing and detuning choose those phasors so they
add forward and partly cancel behind---the same interference calculation as
\cref{sec:arrayfactor}, now with currents found from coupling rather than prescribed
by two feed lines.
\end{description}

\begin{workedbox}[: a crossed Yagi and a fixed dish]
Two equal orthogonal Yagis are fed with equal currents.  If their phase difference
is zero, their fields add along a diagonal and the result is linearly polarized.
Insert a \(90^\circ\) phasing line and the components become
\(E_0\cos\omega t\) and \(E_0\cos(\omega t+90^\circ)\): equal components in
quadrature, hence circular polarization.

Now keep a dish and its feed unchanged but move from \SI{1.2}{\giga\hertz} to
\SI{2.4}{\giga\hertz}.  The wavelength halves, so the aperture formula predicts
\(G_2/G_1=(2.4/1.2)^2=4\), or \SI{6}{\decibel} more gain.  That result is about
the electrically larger aperture, not extra transmitter power.
\end{workedbox}

\begin{mistakebox}
``Resonant,'' ``matched,'' and ``efficient'' answer different questions.  A
loading network can make the feed point look like \SI{50}{\ohm}; it cannot turn
ground loss into radiation.  Check reactance, mismatch, and radiation efficiency
separately before declaring an antenna improved.
\end{mistakebox}
